\section{A second verified transfer: density gaps}
\label{sec:density}

The incoming report \cite{IncomingDensity} addresses McNew's Questions 7.1
and 7.3 \cite{McNew}. Let \(\alpha\) be the supremum of natural densities
of sets of positive integers containing no three-term geometric
progression with rational ratio greater than one, and let
\(\overline\alpha\) use upper density. Define \(\beta,\overline\beta\)
in the same way using only integer ratios.

The verified claims are
\begin{equation}
\overline\alpha-\alpha>0.002836,\qquad
\overline\beta-\beta>0.069218.
\label{eq:density-gaps}
\end{equation}
These are existing incoming results, independently rechecked here.
They are not consequences of RH or GRH. The first strict-gap mechanism
has a short proof that makes their mathematical content assessable.

\subsection{Fixed choices versus choices varying with the cutoff}

Fix \(Y\ge12\). Let \(S\) be the positive integers with all prime factors
at most \(Y\), and put
\[
c=\prod_{p\le Y}(1-p^{-1}),\qquad
v=\prod_{p>Y}(1-p^{-2}).
\]
Write every integer uniquely as \(gh\), with \(h\in S\) and \(g\)
having no prime factor at most \(Y\). Call a set admissible if it
avoids the chosen kind of progression. Define
\[
B=c\sup_{A\subset S\text{ admissible}}\sum_{h\in A}\frac1h,\qquad
U=c\int_1^\infty J(t)\frac{dt}{t^2},
\]
where \(J(t)\) is the largest cardinality of an admissible subset of
\(S\cap[1,t]\). Since the total smooth reciprocal mass is \(1/c\),
one has \(U\le1\).

If an admissible set \(E\) has natural density \(d\), partial summation
gives
\(\sum_{n\le X,\ n\in E}1/n=(d+o(1))\log X\).
Each fixed rough fiber \(A_g=\{h:gh\in E\}\) is admissible. Hence
\[
\sum_{\substack{n\le X\\n\in E}}\frac1n
\le \frac Bc\sum_{\substack{g\le X\\g\ {\rm rough}}}\frac1g
=(B+o(1))\log X,
\]
because roughness here is a finite periodic divisibility condition.
Thus every natural density is at most \(B\).

On the other hand, squarefree rough integers have density \(cv\).
On a progression whose rough parts are squarefree, each outside prime
has three valuations in \(\{0,1\}\) in arithmetic progression, so all
three rough parts agree. This permits independent admissible choices
on each rough fiber. At cutoffs \(X_j\) with \(X_j/X_{j-1}\to\infty\),
assign each new squarefree rough \(g\in(X_{j-1},X_j]\) a maximizing
subset of \(S\cap[1,X_j/g]\), permanently. Earlier fibers can be ignored
for the lower bound at \(X_j\). For fixed \(K\), step-function summation
over \(X_j/K<g\le X_j\) gives the limiting contribution
\[
cv\int_1^K J(t)\frac{dt}{t^2}.
\]
Letting \(K\to\infty\) gives an admissible set of upper density at least
\(vU\).

The difference between \(U\) and \(B\) is strict because no one admissible
set can attain every prefix optimum. Through twelve, the greedy choice
has nine elements,
\[
\{1,2,3,5,6,7,8,10,11\},
\]
whereas
\[
\{1,2,5,6,7,8,9,10,11,12\}
\]
is admissible and has ten. The complete rational-ratio progression list
through twelve is
\[
(1,2,4),\ (1,3,9),\ (2,4,8),\ (3,6,12),\ (4,6,9);
\]
the first four have integer ratios. A set optimal at every earlier
prefix would have to agree with greedy by induction, so there is a
deficit of at least one at some integer cutoff at most twelve.
It persists throughout a unit interval and contributes at least
\(1/(12\cdot13)\) to the integral. Therefore
\[
U-B\ge \frac c{156},\qquad
u-b\ge vU-B\ge \frac c{156}-\sum_{p>Y}p^{-2},
\]
where \(u,b\) are the upper- and natural-density optima for either convention.
The Mertens product estimate and \(\sum_{n>Y}n^{-2}\le1/Y\)
give strict positivity for large \(Y\). More concretely, the exact
certificate at \(Y=2000\) gives a common gap greater than \(0.000404\).

\subsection{The stronger numbers and their scope}

The rational-ratio bound in \eqref{eq:density-gaps} uses all prefix losses
through \(80\). Its exact deficit is \(1553/45360\); the finite covering
optimization and the Euler-product tail bound were independently replayed.
For integer ratios, the certified natural-density upper bound is
\[
\beta\le\frac{5245322990569}{7001316000000}<0.749192.
\]
McNew's Theorem 4.4 gives \(\overline\beta>0.818410\) \cite{McNew}.
Their difference is greater than \(0.069218\). The incoming programs
use integer-weight branch-and-bound, not a floating optimizer.

This answers the strict comparisons posed in those questions, subject
to the written proof and finite certificates. We have not established
historical priority against every later publication. The exact value of
\(\alpha\), Rankin's optimality conjecture, and the comparison of
\(\alpha\) with \(\beta\) are not resolved by this verification.
Other attainment and computability claims in the same incoming archive
are not included in the present certificate claim.
